\documentclass[11pt,a4paper]{article}

\usepackage[margin=2.5cm]{geometry}
\usepackage[T1]{fontenc}
\usepackage[utf8]{inputenc}
\usepackage{hyperref}
\usepackage{xcolor}
\usepackage{listings}

\lstdefinestyle{pythonstyle}{
    language=Python,
    basicstyle=\ttfamily\small,
    keywordstyle=\color{blue!70!black},
    stringstyle=\color{green!40!black},
    commentstyle=\color{gray},
    showstringspaces=false,
    frame=single,
    breaklines=true,
    tabsize=4
}

\title{FEM-PR: istruzioni per lo script del modello}
\author{}
\date{}

\begin{document}

\maketitle

\section{Idea generale}

Il progetto e' organizzato intorno a due classi principali:

\begin{itemize}
    \item \texttt{Domain}: contiene gli oggetti del modello, divisi per categoria
    (\texttt{nodes}, \texttt{materials}, \texttt{sections}, \texttt{beamIntegrations},
    \texttt{geometricTransformations}, \texttt{elements}).
    \item \texttt{Model}: fornisce una sintassi compatta per creare gli oggetti e
    registrarli nel \texttt{Domain}.
\end{itemize}

L'utente normalmente lavora con \texttt{Model}. Per esempio:

\begin{lstlisting}[style=pythonstyle]
from model import Model

m = Model()

m.node(1, 0.0, 0.0)
m.node(2, 0.0, 3.0)
\end{lstlisting}

Ogni metodo di costruzione crea un oggetto, lo aggiunge al \texttt{Domain} e lo
restituisce. Quindi sono equivalenti, come stile d'uso:

\begin{lstlisting}[style=pythonstyle]
m.node(1, 0.0, 0.0)
\end{lstlisting}

oppure:

\begin{lstlisting}[style=pythonstyle]
n1 = m.node(1, 0.0, 0.0)
\end{lstlisting}

Nel primo caso il nodo viene comunque salvato in \texttt{m.nodes}; nel secondo
caso si conserva anche un riferimento locale all'oggetto.

\section{Pattern prototipo e copia}

Alcuni oggetti vengono definiti come prototipi nel \texttt{Domain}. Quando un
altro oggetto ha bisogno di usarli con stato indipendente, viene creata una copia.

\subsection{Materiali}

I materiali definiti dall'utente sono salvati in \texttt{m.materials}. Le sezioni
che hanno bisogno di materiali indipendenti usano copie dei materiali.

\subsection{Sezioni}

Le sezioni sono salvate in \texttt{m.sections}. Gli elementi possono creare copie
delle sezioni, ad esempio una copia per ogni punto di integrazione.

\subsection{Trasformazioni geometriche}

La trasformazione geometrica e' un caso particolare. Nel modello si definisce
solo il tipo e l'identificativo:

\begin{lstlisting}[style=pythonstyle]
m.geomTransf("linear", "linearTransf")
\end{lstlisting}

La trasformazione diventa concreta solo quando viene assegnata a un elemento,
perche' in quel momento sono noti \texttt{nodeI} e \texttt{nodeJ}. Il metodo
\texttt{getGeomTransfCopy(...)} crea quindi una copia inizializzata sui nodi
dell'elemento.

\section{Sintassi dello script del modello}

\subsection{Creazione del modello}

\begin{lstlisting}[style=pythonstyle]
from model import Model

m = Model()
\end{lstlisting}

\subsection{Nodi}

\begin{lstlisting}[style=pythonstyle]
m.node(id, x, y)
\end{lstlisting}

\begin{itemize}
    \item \texttt{id}: identificativo del nodo, intero o stringa.
    \item \texttt{x}, \texttt{y}: coordinate del nodo nel piano.
\end{itemize}

Esempio:

\begin{lstlisting}[style=pythonstyle]
m.node(1, 0.0, 0.0)
m.node(2, 0.0, 3.0)
\end{lstlisting}

\subsection{Materiali}

\subsubsection{\texttt{elastic}}

\begin{lstlisting}[style=pythonstyle]
m.material("elastic", id, E, rho=0.0)
\end{lstlisting}

\begin{itemize}
    \item \texttt{id}: identificativo del materiale.
    \item \texttt{E}: modulo elastico.
    \item \texttt{rho}: densita', opzionale.
\end{itemize}

Esempio:

\begin{lstlisting}[style=pythonstyle]
m.material("elastic", "matElastic", 210000.0)
\end{lstlisting}

\subsubsection{\texttt{steel01}}

\begin{lstlisting}[style=pythonstyle]
m.material("steel01", id, fy, E0, b, a1=0.07, a2=2, a3=0.07, a4=2, rho=0.0)
\end{lstlisting}

\begin{itemize}
    \item \texttt{id}: identificativo del materiale.
    \item \texttt{fy}: tensione di snervamento.
    \item \texttt{E0}: modulo elastico iniziale.
    \item \texttt{b}: rapporto di incrudimento.
    \item \texttt{a1}, \texttt{a2}, \texttt{a3}, \texttt{a4}: parametri opzionali del modello.
    \item \texttt{rho}: densita', opzionale.
\end{itemize}

Esempio:

\begin{lstlisting}[style=pythonstyle]
m.material("steel01", "mat1", 450.0, 210000.0, 0.005)
\end{lstlisting}

\subsection{Sezioni}

\subsubsection{\texttt{elasticSection}}

\begin{lstlisting}[style=pythonstyle]
m.section("elasticSection", id, E, A, I)
\end{lstlisting}

\begin{itemize}
    \item \texttt{id}: identificativo della sezione.
    \item \texttt{E}: modulo elastico.
    \item \texttt{A}: area.
    \item \texttt{I}: momento di inerzia.
\end{itemize}

Esempio:

\begin{lstlisting}[style=pythonstyle]
m.section("elasticSection", "secElastic", 210000.0, 0.02, 8.0e-5)
\end{lstlisting}

\subsubsection{\texttt{aggregator}}

\begin{lstlisting}[style=pythonstyle]
m.section("aggregator", id, axialID, bendingID)
\end{lstlisting}

\begin{itemize}
    \item \texttt{id}: identificativo della sezione.
    \item \texttt{axialID}: identificativo del materiale per la risposta assiale.
    \item \texttt{bendingID}: identificativo del materiale per la risposta flessionale.
\end{itemize}

Esempio:

\begin{lstlisting}[style=pythonstyle]
m.material("steel01", "mat1", 450.0, 210000.0, 0.005)
m.section("aggregator", "sec1", "mat1", "mat1")
\end{lstlisting}

\subsection{Integrazione lungo l'elemento}

\subsubsection{\texttt{lobatto}}

\begin{lstlisting}[style=pythonstyle]
m.beamIntegration("lobatto", id, nIP, sectionID)
\end{lstlisting}

\begin{itemize}
    \item \texttt{id}: identificativo della regola di integrazione.
    \item \texttt{nIP}: numero di punti di integrazione.
    \item \texttt{sectionID}: identificativo della sezione associata.
\end{itemize}

Esempio:

\begin{lstlisting}[style=pythonstyle]
m.beamIntegration("lobatto", "beamInt1", 5, "sec1")
\end{lstlisting}

\subsubsection{\texttt{legendre}}

\begin{lstlisting}[style=pythonstyle]
m.beamIntegration("legendre", id, nIP, sectionID)
\end{lstlisting}

Esempio:

\begin{lstlisting}[style=pythonstyle]
m.beamIntegration("legendre", "beamInt2", 3, "sec1")
\end{lstlisting}

\subsection{Trasformazioni geometriche}

\subsubsection{\texttt{linear}}

\begin{lstlisting}[style=pythonstyle]
m.geomTransf("linear", id)
\end{lstlisting}

\begin{itemize}
    \item \texttt{id}: identificativo della trasformazione geometrica.
\end{itemize}

Esempio:

\begin{lstlisting}[style=pythonstyle]
m.geomTransf("linear", "linearTransf")
\end{lstlisting}

\subsection{Elementi}

\subsubsection{\texttt{dispBeamColumn}}

\begin{lstlisting}[style=pythonstyle]
m.element("dispBeamColumn", id, nodeI_ID, nodeJ_ID, geomTransfID, beamIntID)
\end{lstlisting}

\begin{itemize}
    \item \texttt{id}: identificativo dell'elemento.
    \item \texttt{nodeI\_ID}: identificativo del nodo iniziale.
    \item \texttt{nodeJ\_ID}: identificativo del nodo finale.
    \item \texttt{geomTransfID}: identificativo della trasformazione geometrica.
    \item \texttt{beamIntID}: identificativo della regola di integrazione.
\end{itemize}

Esempio:

\begin{lstlisting}[style=pythonstyle]
m.element("dispBeamColumn", "dispBC", 1, 2, "linearTransf", "beamInt1")
\end{lstlisting}

Quando l'elemento viene creato, il modello:

\begin{enumerate}
    \item recupera i due nodi dal \texttt{Domain};
    \item crea una copia della trasformazione geometrica inizializzata sui due nodi;
    \item crea una copia della \texttt{beamIntegration};
    \item costruisce l'elemento \texttt{DispBeamColumn};
    \item aggiunge l'elemento a \texttt{m.elements}.
\end{enumerate}

All'interno di \texttt{DispBeamColumn}, il metodo \texttt{initialize()} crea una
copia della sezione per ogni punto di integrazione e inizializza vettori e matrici
dell'elemento.

\section{Esempio completo}

\begin{lstlisting}[style=pythonstyle]
from model import Model

m = Model()

m.node(1, 0.0, 0.0)
m.node(2, 0.0, 3.0)

m.material("steel01", "mat1", 450.0, 210000.0, 0.005)
m.section("aggregator", "sec1", "mat1", "mat1")

m.beamIntegration("lobatto", "beamInt1", 5, "sec1")
m.geomTransf("linear", "linearTransf")

ele = m.element("dispBeamColumn", "dispBC", 1, 2, "linearTransf", "beamInt1")

print(m.nodes)
print(m.materials)
print(m.sections)
print(m.beamIntegrations)
print(m.geometricTransformations)
print(m.elements)
\end{lstlisting}

\section{Metodi espliciti disponibili}

Oltre alla sintassi compatta, \texttt{Model} espone anche metodi espliciti:

\begin{lstlisting}[style=pythonstyle]
m.addElasticMaterial(id, E, rho=0.0)
m.addSteel01Material(id, fy, E0, b, a1=0.07, a2=2, a3=0.07, a4=2, rho=0.0)

m.addElasticSection(id, E, A, I)
m.addAggregator(id, axialID, bendingID)

m.addGaussLobatto(id, nIP, sectionID)
m.addGaussLegendre(id, nIP, sectionID)

m.addLinearGeomTransf(id)

m.addDispBeamColumn(id, nodeI_ID, nodeJ_ID, geomTransfID, beamIntID)
\end{lstlisting}

Questi metodi sono piu' espliciti e possono essere comodi durante debug,
documentazione o uso con suggerimenti dell'IDE.

\end{document}
